\begin{abstract}

Cellular core networks (CNs) are critical infrastructure, yet their internal security model has historically relied on physical isolation: interfaces between core components often operate within an assumed trust zone. As CNs transition to cloud-native deployments, this assumption weakens, expanding the attack surface and enabling external adversaries to reach previously internal interfaces.
From a root-cause analysis of security flaws reported in GitHub issues for open-source CN implementations, we found a recurring pattern of \textit{blind trust} among CN components. Components may omit syntactic validation, fail to enforce semantic invariants, or allocate resources without checking availability. Once internal interfaces become reachable, these weaknesses can lead to severe impacts such as denial of service and session hijacking. We call these vulnerabilities \textit{implicit trust errors} (\textit{iTrue}).

To detect iTrues and understand their security impacts, we designed \textit{iFinder}, an LLM-driven multi-agent system that summarizes known flaws, distills them into detection patterns, and applies them to discover new iTrues in CN implementations. 
To suppress hallucinations produced by large language models (LLMs), we built an innovative strategy that cross-checks both 3GPP specifications and CN code to capture existing protection missed by the agents. Further, we developed a technique that uses LLMs to generate proof-of-concept (PoC) exploits for potential iTrues and iteratively refine the PoCs by automatically executing them against CN implementations and analyzing results. 
Running iFinder on seven prominent open-source CN implementations, we discovered 84 previously unknown vulnerabilities. Among them, 83 have already been confirmed and 81 have been assigned CVEs. 
Importantly, a session-hijacking flaw has been confirmed on real-world commercial 5G core networks. Our findings highlight the pervasiveness of iTrue risks across the cellular core networks and the urgent need for elevated protection within the original trust domains. We plan to make iFinder publicly available to help enhance the security of cellular core networks. 
\end{abstract}
